\documentclass[11pt,a4paper]{article}
\usepackage{curso}
\begin{document}
\portada{Lección}{3.5}{Diseño de protocolos y aceleración de reloj}

\noindent\textbf{Unidad 3.}\quad\textbf{Objetivos:} carga súbita frente a gradual; aceleración de reloj y sus leyes de escala.
\textbf{Prerrequisitos:} 3.1--3.4.

\section{Cargar un problema}
Llevar el hardware desde un estado inicial a la energía $E_{\text{problema}}$ es un protocolo $\lambda(t)$. Ejemplo: $\kappa:1\to4$
($\kB T=\mu=1$, $\Delta F=0{,}693$). Carga súbita desde equilibrio: $W=\frac12(\kappa_f-\kappa_i)\langle x^2\rangle_i=\frac{\kB T}2(\frac{\kappa_f}{\kappa_i}-1)=1{,}5$.
Rampa lineal de duración $\tau$ \evS{} ($4\times10^4$ trayectorias):
\begin{center}\begin{tabular}{lccccccc}\toprule
$\tau$ & súbito & 0,25 & 0,5 & 1 & 2 & 4 & 8\\\midrule
$\langle W\rangle$ & 1,51 & 1,23 & 1,08 & 0,94 & 0,83 & 0,77 & 0,73\\\bottomrule\end{tabular}\end{center}
Para $\tau$ grande, $W_{\text{dis}}\approx a/\tau$ ($a\approx0{,}3$), coherente con el límite de velocidad de 3.4.

\section{Aceleración de reloj}
\begin{resultado}
Si $dx=-\mu\lambda V'(x)dt+\sqrt{2\mu\lambda\kB T}\,dW$ y $s=\lambda t$, entonces $dx=-\mu V'(x)ds+\sqrt{2\mu\kB T}\,d\tilde W(s)$, porque
$dW(t)=\lambda^{-1/2}d\tilde W(s)$ en ley. Las trayectorias son las mismas, recorridas $\lambda$ veces más rápido \cite{whitelam}.
\end{resultado}
Escalar $V$ y la intensidad del ruido equivale a subir la temperatura por $\lambda$: la distribución estacionaria $\propto e^{-\lambda V/(\lambda\kB T)}$
no cambia. Consecuencias: energías y calor $\times\lambda$; duración $\div\lambda$; potencia $\times\lambda^2$; $\Sigma=Q/(\lambda\kB T)$ y $\eta$ invariantes.
\begin{ejemplo}
Un cómputo que disipa $10\,\kB T$ en $1\,\mu$s, con $\lambda=10$: 100 ns, $100\,\kB T$, potencia ×100, $\Sigma=10$ igual.
\end{ejemplo}

\begin{aplicacion}
Laboratorio L4.4: verificar invariancia de $\Sigma$ y $\eta$ y el escalado de calor y potencia; reto: tiempo de carga óptimo en una red cuártica.
\end{aplicacion}

\begin{error}
\begin{itemize}[nosep]
\item Escalar $V$ sin escalar el ruido cambia la distribución (temperatura efectiva menor).
\item $\Sigma$ invariante no significa gratis: calor y potencia crecen. \quad El ancho de banda del hardware limita $\lambda$.
\end{itemize}
\end{error}

\begin{thebibliography}{9}\small
\bibitem[42]{whitelam} S. Whitelam, The thermodynamic freedom of a thermodynamic computer, 2026. arXiv:2608.27938.
\bibitem[43]{wc} S. Whitelam y C. Casert, Thermodynamic computing out of equilibrium, 2024.
\end{thebibliography}
\end{document}
